\begin{proof}[Proof of \cref{obj:thm:observable-upper}]
\begin{enumerate}
\item Fix public class constants \(K\), \(\beta\in(0,1]\), and \(\kappa\in[0,2]\). The Fourier and polynomial comparison-score conclusions of \cref{obj:lem:dictionary-score-rate} supply positive constants \(C_{\mathrm F},C_{\mathrm M}\) and integer thresholds \(n_{\mathrm F},n_{\mathrm M}\), depending on \(K,\beta,\kappa\), such that
\[
\begin{aligned}
n\ge n_{\mathrm F},\quad \sigma\le L_n^{-1/2}
&\ \Longrightarrow\
\exists k\ge0:\ 
\text{the Fourier pair with index }k\text{ belongs to }\mathcal D_{n,\sigma},
\quad A_q(K)\le C_{\mathrm F}\rho_{n,\sigma},\\
n\ge n_{\mathrm M},\quad L_n^{-1/2}<\sigma
&\ \Longrightarrow\
\exists j\ge0:\ 
\text{the polynomial pair with index }2^j+1\text{ belongs to }\mathcal D_{n,\sigma},
\quad A_q(K)\le C_{\mathrm M}\rho_{n,\sigma},
\end{aligned}
\]
for every \(\sigma\in[0,1/4]\). These inequalities concern public scores and rates determined by \(\beta,\kappa,n,\sigma\), so the constants and thresholds can be fixed uniformly over the potential-path spaces. Set
\[
C=\max\{C_{\mathrm F},C_{\mathrm M}\},
\qquad
n_0=\max\{3,n_{\mathrm F},n_{\mathrm M}\}.
\]
Then \(C>0\). % lean: CausalSmith.Stat.NoisydoseWeakdesignTransition.fourier_dictionary_witness; CausalSmith.Stat.NoisydoseWeakdesignTransition.inverseheat_dictionary_witness

\item Fix a potential-path space, an integer \(n\ge n_0\), and \(\sigma\in[0,1/4]\). In particular,
\[
n\ge3,\qquad n\ge n_{\mathrm F},\qquad n\ge n_{\mathrm M},
\qquad L_n=\log(en)\ge0.
\]
The three expressions defining \(b_{n,\sigma}\) in \cref{obj:def:frontier-rate} are nonnegative: the direct expression is a power of a positive sample size, the Fourier expression has nonnegative numerator and positive denominator, and the polynomial expression has numerator \(\log(e+\sigma^2L_n)\ge0\) and denominator \(L_n>0\). Consequently,
\[
b_{n,\sigma}\ge0,
\qquad
\rho_{n,\sigma}=b_{n,\sigma}^{\beta}\ge0.
\]
This permits enlargement of either comparison constant to \(C\). % lean: CausalSmith.Stat.NoisydoseWeakdesignTransition.observable_upper

\item Suppose \(\sigma\le L_n^{-1/2}\). Choose the Fourier comparison pair from the first witness bound. Its dictionary membership and \(n>0\) allow application of the admissibility conclusion of \cref{obj:lem:dictionary-score-rate}. Thus its functions are measurable and satisfy
\[
\begin{gathered}
q(t)\ge0\quad(t\in[-1/2,1/2]),
\qquad 0<I_\kappa(q)<\infty,\\
E_Z|\ell_q(t+\sigma Z)|<\infty,
\qquad E_Z[\ell_q(t+\sigma Z)]=q(t)
\quad(t\in[-1/2,1/2]),
\qquad V_q<\infty.
\end{gathered}
\]
Moreover,
\[
A_q(K)\le C_{\mathrm F}\rho_{n,\sigma}
\le C\rho_{n,\sigma},
\]
where the last inequality follows from \(C_{\mathrm F}\le C\) and \(\rho_{n,\sigma}\ge0\). This proves the Fourier comparison guarantee, including equality at the cutoff. % lean: CausalSmith.Stat.NoisydoseWeakdesignTransition.observable_upper.hFourier

\item Suppose \(L_n^{-1/2}<\sigma\). Choose the polynomial comparison pair from the second witness bound, with
\[
m=2^j+1,
\qquad
(q,\ell_q)=(q_m^{\mathrm M},\ell_m^{\mathrm M}).
\]
Its dictionary membership again gives, by \cref{obj:lem:dictionary-score-rate}, measurability and
\[
\begin{gathered}
q(t)\ge0\quad(t\in[-1/2,1/2]),
\qquad 0<I_\kappa(q)<\infty,\\
E_Z|\ell_q(t+\sigma Z)|<\infty,
\qquad E_Z[\ell_q(t+\sigma Z)]=q(t)
\quad(t\in[-1/2,1/2]),
\qquad V_q<\infty.
\end{gathered}
\]
Its score satisfies
\[
A_q(K)\le C_{\mathrm M}\rho_{n,\sigma}
\le C\rho_{n,\sigma}.
\]
This proves the polynomial comparison guarantee. % lean: CausalSmith.Stat.NoisydoseWeakdesignTransition.observable_upper.hPoly

\item Let \(\widehat q\) be the selected latent weight in \cref{obj:def:total-estimator}. The fixed-order minimization gives
\[
(\widehat q,\ell_{\widehat q})\in\mathcal D_{n,\sigma},
\qquad
A_{\widehat q}(K)\le A_q(K)
\quad\text{for every }(q,\ell_q)\in\mathcal D_{n,\sigma}.
\]
If \(\sigma\le L_n^{-1/2}\), compare with the Fourier witness; otherwise \(L_n^{-1/2}<\sigma\), and compare with the polynomial witness. In either case,
\[
A_{\widehat q}(K)\le A_q(K)\le C\rho_{n,\sigma}.
\]
% lean: CausalSmith.Stat.NoisydoseWeakdesignTransition.selectedTag_minimizes; CausalSmith.Stat.NoisydoseWeakdesignTransition.observable_upper.hscore

\item Fix \(P\in\mathcal P_{K,\beta,\kappa,\sigma}\) satisfying \cref{obj:ass:iid}. Applied to the selected dictionary pair, \cref{obj:lem:dictionary-score-rate} supplies all its weight-admissibility conditions, finite extended second moment, and the law-specific integrability
\[
E_P\!\left[|\ell_{\widehat q}(W-a_0)|\right]<\infty.
\]
Hence \cref{obj:lem:public-error-certificate} yields
\[
E_{Q_P^n}\!\left[
|\widehat\theta_{\widehat q}-\theta(P)|
\right]\le A_{\widehat q}(K).
\]
For each \(r\in\{0,1,2\}\), the sample index \(i=r\) exists because \(n\ge3\), and \(i\bmod3=r\). Therefore
\[
r\in\mathcal I_r,\qquad
\mathcal I_r\ne\varnothing\quad(r=0,1,2),
\qquad
\widehat\theta=\widehat\theta_{\widehat q}.
\]
Using the total-estimator formula in \cref{obj:def:total-estimator} and the selected-score bound,
\[
E_{Q_P^n}\!\left[|\widehat\theta-\theta(P)|\right]
=
E_{Q_P^n}\!\left[
|\widehat\theta_{\widehat q}-\theta(P)|
\right]
\le A_{\widehat q}(K)
\le C\rho_{n,\sigma}.
\]
The choices of \(C,n_0\) apply to every permitted space, law, and public error scale, as required. % lean: CausalSmith.Stat.NoisydoseWeakdesignTransition.public_error_certificate; CausalSmith.Stat.NoisydoseWeakdesignTransition.blocksNonempty_of_three; CausalSmith.Stat.NoisydoseWeakdesignTransition.observable_upper
\end{enumerate}
\end{proof}
